% Study plan 14:30-19:00, subchapter by subchapter. GENERATED by gen_plan.py: edit there.
\providecommand{\G}[1]{\S\ref{#1}, p.\,\pageref{#1}}
\providecommand{\tick}{$\square$}
\newtcolorbox{chapplan}[1]{colback=white,colframe=navy,title={#1},left=2mm,right=2mm}
\newcommand{\sub}[5]{\par\smallskip\noindent\textbf{\color{navy}#1}\ \ {\small(#2 min)}\ \ \textbf{#4} \hfill{\small\color{grey}\G{#3}}\par
 \noindent\hspace*{4mm}\parbox{\dimexpr\linewidth-4mm}{\textbf{Do:} #5}\par}
\newcommand{\gate}[1]{\noindent\hspace*{4mm}\parbox{\dimexpr\linewidth-4mm}{\tick\ \textbf{Move on when you can} #1}\par\smallskip}
\newcommand{\pausebox}[2]{\begin{tcolorbox}[colback=green!6,colframe=green!60!black,title={Break \quad #1}]#2\end{tcolorbox}}

\section{Your study plan: 14:30 to 19:00, subchapter by subchapter}\label{sec:plan}

\subsection*{How it works}
You go through the guide \textbf{one subchapter at a time}, in order. Each subchapter has a time slot, a \textbf{Do} list and a
\textbf{Move on when you can\dots} test. Only move to the next subchapter when you can pass its test \emph{without looking}.

\begin{methodbox}[title={The loop for every subchapter}]
\begin{enumerate}
\item \textbf{Read} the definition and method boxes slowly. For every step ask ``\emph{why} is this allowed?''
\item \textbf{Cover} each worked example and do it yourself on paper; uncover one line only when stuck.
\item \textbf{Explain} the main idea aloud in one or two sentences, as if to a friend.
\item \textbf{Test}: close the guide and do the ``Move on when you can'' task. Pass $\Rightarrow$ tick the box and move on.
\end{enumerate}
\end{methodbox}
\begin{warnbox}[title={The stuck rule (keeps you on time without skipping understanding)}]
If you still cannot pass a test \textbf{3 minutes after its slot ends}, write the subchapter number on a ``come back'' list, put a star next to it, and move on.
Later subchapters often make earlier ones click. The final slot (18:50--19:00) is reserved for your starred subchapters.
Never just reread: always \emph{do} something (write, derive, explain aloud).
\end{warnbox}
Have ready: the guide, plain paper, pen, calculator, a ``come back'' list, phone in another room.

\subsection*{Timetable at a glance}
\begin{center}\renewcommand{\arraystretch}{1.25}\small
\begin{tabular}{@{}p{2.3cm}p{1cm}p{7.5cm}p{3.4cm}@{}}\toprule
\textbf{Time} & \textbf{Min} & \textbf{Part} & \textbf{Guide}\\\midrule
14:30--14:35 & 5 & Overview & \G{sec:map}\\
14:35--15:11 & 36 & Chapter 1: Describing linear systems & \G{ch1}\\
15:11--15:16 & 5 & Break & \\
15:16--16:22 & 66 & Chapter 2: Matrix theory & \G{ch2}\\
16:22--16:32 & 10 & Break & \\
16:32--17:20 & 48 & Chapter 3: Laplace transforms & \G{ch3}\\
17:20--17:25 & 5 & Break & \\
17:25--17:55 & 30 & Chapter 4: Frequency response and stability & \G{ch4}\\
17:55--18:00 & 5 & Break & \\
18:00--18:50 & 50 & Chapter 5: Feedback control & \G{ch5}\\
18:50--19:00 & 10 & Come-back list and recall & \G{sec:proofs}\\
\bottomrule\end{tabular}\end{center}
Each chapter is split into its subchapters on the following pages, with the minutes chosen by how hard and how examinable each one is.
\subsection*{Subchapter by subchapter}
\begin{chapplan}{Start \quad 14:30--14:35 \quad The whole module on one page}
\sub{14:30--14:35}{5}{sec:map}{The big picture}{Read the chapter table and the chain diagram. Then, closed book, write the chain ODE $\to$ diagram $\to$ $(A,B,C,D)$ $\to$ $T(s)$ $\to$ stability $\to$ feedback.}
\gate{write $\dot X=AX+BU$, $Y=CX+DU$ and $T(s)=D+C(sI-A)^{-1}B$ from memory.}
\end{chapplan}
\begin{chapplan}{Chapter 1: Describing linear systems \quad 14:35--15:11 (36 min)}
\sub{14:35--14:39}{4}{ss:1-1}{Linear operators and the five building blocks}{Read the definition box. Say what each block does to a signal $u$. Do the exam tip: show $Lf=tf''+3\int_0^tf$ is linear.}
\gate{say what a tap, summing junction, integrator, differentiator and amplifier do, and check linearity of a given $L$.}
\sub{14:39--14:44}{5}{ss:1-2}{Block diagrams: open loop and feedback}{Cover and redo Black's amplifier ($y=u+\theta y\Rightarrow y=u/(1-\theta)$) and the cyclist ($v=-\frac km\int v+\frac1m\int u$).}
\gate{tell open loop from feedback, and write the equation at a summing junction.}
\sub{14:44--14:54}{10}{ss:1-3}{How to turn an equation into a block diagram}{Learn the 5-step method box. Cover and redo A1.1(i), A1.1(ii) (integrate twice, solve for $y$, nest the integrals) and A1.2.}
\gate{turn $y''+4y=u'$ (or similar) into a diagram with \emph{no differentiators}, closed book.}
\sub{14:54--15:00}{6}{ss:1-4}{Reduction of order and the matrix form}{Learn the 4-step method. Cover and redo A1.3 ($m\ddot x+kx=u$) and the damped oscillator.}
\gate{write $A,B,C,D$ for any $n$-th order constant-coefficient ODE (remember to make it monic first).}
\sub{15:00--15:05}{5}{ss:1-5}{The companion matrix}{Learn the sign/order rule (last row $-a_0,\dots,-a_{n-1}$). Redo A2.2. Read the induction proof once.}
\gate{write a matrix with characteristic polynomial $s^3-2s^2+s+7$ in 30 seconds.}
\sub{15:05--15:11}{6}{ss:1-6}{Exam-style questions for Chapter 1}{Do Q1.1 and Q1.2 closed book, then mark against the solutions. Q1.3 if time.}
\gate{get Q1.1 and Q1.2 fully right.}
\end{chapplan}
\pausebox{15:11--15:16}{Stand up, get water, no phone.}
\begin{chapplan}{Chapter 2: Matrix theory \quad 15:16--16:22 (66 min)}
\noindent\textbf{Warm-up recall (inside the first slot, 1 min, closed book):} $A,B,C,D$ for $m\ddot x+kx=u$; the companion matrix of $s^3+2s+5$.\par
\sub{15:16--15:20}{4}{ss:2-1}{Determinants, cofactors, characteristic polynomial}{Read the box. Note: $\chi_A(s)=\det(sI-A)$, adjugate $=$ \emph{transpose} of cofactors, $2\times2$: $s^2-\tr A\,s+\det A$.}
\gate{write $\chi_A$ for a $2\times2$ matrix instantly and explain what the adjugate is.}
\sub{15:20--15:27}{7}{ss:2-2}{Computing the resolvent $(sI-A)^{-1}$}{Learn the 5-step method. Cover and redo A1.4 (upper triangular, so many cofactors are 0).}
\gate{compute $(sI-A)^{-1}$ for a $2\times2$ matrix in under 3 minutes and check it.}
\sub{15:27--15:30}{3}{ss:2-3}{Norms and Cauchy--Schwarz}{Read the results box and the Cauchy--Schwarz proof once (discriminant $\le0$ idea).}
\gate{state $\|AB\|\le\|A\|\|B\|$ and the one-line idea of the Cauchy--Schwarz proof.}
\sub{15:30--15:34}{4}{ss:2-4}{Rank and the matrices $Q$, $R$ (AE2)}{Redo A2.3(i)--(ii). Read the ``why rank matters'' box (rank drop $\leftrightarrow$ pole cancels).}
\gate{prove $Q^\top=R$ and find $R$ when $C^\top$ is an eigenvector of $A^\top$.}
\sub{15:34--15:40}{6}{ss:2-5}{Eigenvalues, diagonalisation, Cayley--Hamilton}{Read the results box and the $AS=S\Delta$ proof. Cover and redo the Cayley--Hamilton example ($A^{-1}$ and $A^3$).}
\gate{explain why $g(A)=Sg(\Delta)S^{-1}$ and use Cayley--Hamilton to invert a $2\times2$.}
\sub{15:40--15:53}{13}{ss:2-6}{The matrix exponential}{Learn the 3-way method box. Cover and redo: distinct eigenvalues, repeated eigenvalue ($e^{\lambda t}(I+t(A-\lambda I))$), $\pm i$ (rotation). Check $t=0$ every time. Read A3.3.}
\gate{compute $\exp(tA)$ for a $2\times2$ with distinct \emph{and} with repeated eigenvalues, and check it.}
\sub{15:53--16:00}{7}{ss:2-7}{The solution formula for $\dot X=AX+BU$}{Say the integrating-factor proof aloud. Cover and redo the step-response example; check $y(\infty)=T(0)$.}
\gate{write $X(t)=\exp(tA)X_0+\int_0^t\exp((t-s)A)BU(s)\,ds$ and explain each term (CF and PI).}
\sub{16:00--16:02}{2}{ss:2-8}{SISO and MIMO $(A,B,C,D)$}{Read the box: names and shapes of $A,B,C,D$.}
\gate{give the shapes of $A,B,C,D$ for $k$ inputs, $n$ states and $m$ outputs.}
\sub{16:02--16:09}{7}{ss:2-9}{Rational functions and the transfer function}{Learn proper/strictly proper, poles/zeros, $T=D+C(sI-A)^{-1}B$. Redo the damped-oscillator example; read A2.1 and the proof of Prop 0.20.}
\gate{compute $T(s)$ for a $2\times2$ system and say why every pole is an eigenvalue (but not conversely).}
\sub{16:09--16:17}{8}{ss:2-10}{Realisation: from $T(s)$ back to $(A,B,C,D)$}{Learn the 4-step method. Cover and redo the frame 79 example and the proper example with long division.}
\gate{realise any proper $T(s)$ (including the $D\neq0$ case) and check it by the determinant formula.}
\sub{16:17--16:22}{5}{ss:2-11}{Exam-style questions for Chapter 2}{Do Q2.1 closed book and mark it. Q2.2 if time.}
\gate{get Q2.1 fully right.}
\end{chapplan}
\pausebox{16:22--16:32}{Longer break: eat something. You are halfway.}
\begin{chapplan}{Chapter 3: Laplace transforms \quad 16:32--17:20 (48 min)}
\noindent\textbf{Warm-up recall (inside the first slot, 1 min, closed book):} the solution formula for $\dot X=AX+BU$; $\exp(tA)$ for a double eigenvalue; $T(s)$.\par
\sub{16:32--16:36}{4}{ss:3-1}{Definition and the table}{Read the definition. Cover the table and write it from memory twice.}
\gate{write the whole table from memory.}
\sub{16:36--16:40}{4}{ss:3-2}{Properties}{Learn the 6 properties; say the derivative-rule proof aloud (integration by parts).}
\gate{state $\Lap(f')$, $\Lap(f'')$, the shift rule and the $t f(t)$ rule.}
\sub{16:40--16:45}{5}{ss:3-3}{Laplace transforms from the definition}{Cover and redo A3.1 ($\cos2\omega t$ via $e^{\pm2i\omega t}$, then $\sin^2\omega t$). Read the $x^2$ example.}
\gate{find $\Lap(e^{at})$ and $\Lap(\cos at)$ from the definition, stating the condition on $s$.}
\sub{16:45--16:50}{5}{ss:3-4}{Partial fractions and inverting}{Learn the 5-step invert method and Lerch. Cover and redo the repeated-pole example.}
\gate{invert a function with a repeated pole and one with an irreducible quadratic (complete the square).}
\sub{16:50--16:59}{9}{ss:3-5}{Solving initial value problems}{Learn the 5-step IVP method. Cover and redo A3.2 and the second-order step example.}
\gate{solve a second-order IVP by Laplace from start to finish, including the check.}
\sub{16:59--17:04}{5}{ss:3-6}{Convolution and integral equations}{Learn $\Lap(f*g)=\hat f\hat g$; read the proof. Cover and redo $y=t+\int_0^t\sin(t-u)y(u)\,du$.}
\gate{spot a convolution in an integral equation and solve it.}
\sub{17:04--17:09}{5}{ss:3-7}{Laplace transform of $(A,B,C,D)$}{Derive $\hat Y=T\hat U+C(sI-A)^{-1}X(0)$ yourself. Read the complex-eigenvalue $\exp(tA)$ example.}
\gate{explain why $\hat Y=T(s)\hat U$ when $X(0)=0$, and get $\exp(tA)$ as $\Lap^{-1}[(sI-A)^{-1}]$.}
\sub{17:09--17:15}{6}{ss:3-8}{Jordan canonical form}{Learn: algebraic mult.\ = total block size; geometric mult.\ = number of blocks; minimal polynomial uses the largest block. Read the corrected frame 119 example.}
\gate{read off multiplicities and the minimal polynomial from a Jordan form, and find $\exp(tJ_2(\lambda))$.}
\sub{17:15--17:20}{5}{ss:3-9}{Exam-style questions for Chapter 3}{Do Q3.2 and Q3.4 closed book and mark them.}
\gate{get Q3.2 and Q3.4 fully right.}
\end{chapplan}
\pausebox{17:20--17:25}{Short break.}
\begin{chapplan}{Chapter 4: Frequency response and stability \quad 17:25--17:55 (30 min)}
\noindent\textbf{Warm-up recall (inside the first slot, 1 min, closed book):} $\Lap(f')$, $\Lap(f'')$, $\Lap(\sin at)$, $\Lap(f*g)$.\par
\sub{17:25--17:27}{2}{ss:4-1}{Transfer functions and realisation}{Read the box (Theorems 0.38 and 0.46).}
\gate{state what can be realised by a SISO, and by a BIBO-stable SISO.}
\sub{17:27--17:35}{8}{ss:4-2}{Frequency response, gain and phase}{Learn the method and the quadrant warning. Cover and redo $\frac1{s+1}$ including its Nyquist circle; read the $\frac1{as^2+bs+1}$ example.}
\gate{find $\Gamma(\omega)$, $\phi(\omega)$ for a simple $T$ and sketch its Nyquist plot.}
\sub{17:35--17:41}{6}{ss:4-3}{Stability: the three cases}{Learn the three-row table and the method. Go through the five worked examples and say why each is classified that way.}
\gate{classify $\dot X=AX$ as exponentially stable, marginal or unstable, including the Jordan-block case.}
\sub{17:41--17:44}{3}{ss:4-4}{The damped harmonic oscillator}{Read the table. Connect it to ``$s^2+\beta s+\gamma$ stable iff $\beta,\gamma>0$''.}
\gate{say for which $\beta,\gamma$ the oscillator is stable, marginal or unstable.}
\sub{17:44--17:47}{3}{ss:4-5}{Resonance}{Read the undamped oscillator example. Understand the double pole $\Rightarrow$ $t\sin\nu t$.}
\gate{explain in two sentences why marginally stable is \emph{not} BIBO stable.}
\sub{17:47--17:52}{5}{ss:4-6}{BIBO stability theorems}{Say the proof of Theorem 0.42 aloud (bound $\|\exp(tA)\|\le Me^{-\delta t}$, integrate). Read the outline of Theorem 0.45.}
\gate{write the proof of Theorem 0.42 closed book.}
\sub{17:52--17:55}{3}{ss:4-7}{Exam-style questions for Chapter 4}{Do Q4.1 closed book and mark it.}
\gate{get Q4.1 fully right.}
\end{chapplan}
\pausebox{17:55--18:00}{Short break. Last chapter next.}
\begin{chapplan}{Chapter 5: Feedback control \quad 18:00--18:50 (50 min)}
\noindent\textbf{Warm-up recall (inside the first slot, 1 min, closed book):} the stability table for $\dot X=AX$; the definition of BIBO.\par
\sub{18:00--18:05}{5}{ss:5-1}{Stable polynomials}{Learn Prop 0.43 and its proof, the counterexample $(s+1)(s^2+1)$, and the degree 2 and 3 tests.}
\gate{decide whether a polynomial of degree $\le3$ is stable, and prove Prop 0.43.}
\sub{18:05--18:08}{3}{ss:5-2}{The ring $\mathcal S$}{Learn the two conditions (proper, poles in open LHP) and why $\mathcal S$ is not a field.}
\gate{say whether $\frac{s^2}{s+1}$, $\frac{1}{s^2+1}$, $\frac{s}{(s+1)^2}$ are in $\mathcal S$, with reasons.}
\sub{18:08--18:13}{5}{ss:5-3}{Nyquist's criterion}{Learn the statement and $N=Z-P$. Cover and redo the $\frac{k}{(s+1)^3}$ example (crossing at $-k/8$).}
\gate{find the gain range for stability from where the Nyquist plot crosses the negative real axis.}
\sub{18:13--18:19}{6}{ss:5-4}{The simple feedback loop and PID}{Derive $T=G/(1+GK)$ from the diagram. Cover and redo the PI example and the proportional example.}
\gate{find a PI or PD controller that stabilises a simple plant, stating your sign convention.}
\sub{18:19--18:23}{4}{ss:5-5}{Euclid, B\'ezout and matrix completion}{Learn the algorithm. Cover and redo worked examples (a) and (b); read (c) and (d).}
\gate{find $M,N$ with $PM+QN=1$ for two coprime polynomials of low degree.}
\sub{18:23--18:32}{9}{ss:5-6}{Coprime factorisation and Youla}{Learn the 4-step algorithm and the four formulas $QY,PY,XQ,PX$. Cover and redo worked example A ($\frac1{s-1}$); read B.}
\gate{do the whole algorithm for $G=\frac1{s-a}$ and check all four closed-loop functions.}
\sub{18:32--18:36}{4}{ss:5-7}{Well-posedness and internal stability}{Read Lemma 0.51 and the ``trap'' example.}
\gate{explain why all four functions must be in $\mathcal S$, using the trap example.}
\sub{18:36--18:40}{4}{ss:5-8}{Positive definite matrices}{Learn the three equivalent conditions. Cover and redo A3.4 (i)--(v) and the minors example.}
\gate{test a $3\times3$ matrix by leading minors and prove A3.4(iii).}
\sub{18:40--18:44}{4}{ss:5-9}{Lyapunov's criterion and equation}{Say the proof aloud ($V=\langle KX,X\rangle$ is decreasing). Redo the $K=I$ example; read the Lyapunov-equation example.}
\gate{use $K=I$ to show solutions are bounded, and write the proof.}
\sub{18:44--18:46}{2}{ss:5-10}{Sylvester's equation}{Read the result and the $1\times1$ intuition $a\neq-b$.}
\gate{state when $AY+YB=-C$ has a unique solution for every $C$.}
\sub{18:46--18:50}{4}{ss:5-11}{Exam-style questions for Chapter 5}{Do Q5.1 and Q5.3 closed book and mark them.}
\gate{get Q5.1 and Q5.3 fully right.}
\end{chapplan}
\begin{chapplan}{Finish \quad 18:50--19:00 \quad Come-back list and recall}
\sub{18:50--18:57}{7}{sec:proofs}{Your starred subchapters}{Go back to each subchapter on your ``come back'' list and retry its test. If the list is empty, do mock question M2 (\G{sec:mock}).}
\gate{pass the test of at least one starred subchapter.}
\sub{18:57--19:00}{3}{sec:mock}{Formula sheet from memory}{Cover the formula sheet (last page) and write as many lines as you can; mark what you missed.}
\gate{write at least 15 of the 20 formula-sheet lines.}
\end{chapplan}
\subsection*{Come-back list}
\begin{center}\renewcommand{\arraystretch}{1.6}
\begin{tabular}{@{}p{2.5cm}p{8cm}p{3.5cm}@{}}\toprule
\textbf{Subchapter} & \textbf{What I could not do} & \textbf{Fixed? (time)}\\\midrule
 & & \\ & & \\ & & \\ & & \\ & & \\ & & \\\bottomrule\end{tabular}\end{center}
\begin{exambox}[title={Next session}]
Start with 10 minutes writing the formula sheet from memory, then your come-back list, then the proofs list (\G{sec:proofs}) and the full mock paper (\G{sec:mock}).
\end{exambox}
\newpage
